\section{Backend – Search}

\subsection{Panoramica}

Data l'importanza centrale della funzionalità di ricerca, tale modulo verrà trattato in modo approfondito, con una descrizione dettagliata dei componenti, 
dei flussi di dati e delle logiche di trasformazione, al fine di fornire una comprensione completa del suo funzionamento interno.
Il modulo espone al frontend, tramite IPC di Electron, le funzionalità di interrogazione del pacchetto DiP indicizzato.
Il punto di ingresso è \textit{SearchIpcAdapter}, che registra su \textit{IpcMain} i canali di comunicazione e delega l'elaborazione ai rispettivi Use Case.

Il modulo supporta quattro modalità di ricerca principali:

\begin{itemize}
    \item \textbf{Ricerca per classe documentale} – ricerca le \textit{DocumentClass} per nome tramite \textit{SearchDocumentalClassUC}.
    \item \textbf{Ricerca per processo} – ricerca i \textit{Process} per UUID tramite \textit{SearchProcessUC}.
    \item \textbf{Ricerca strutturata per metadati} – ricerca i \textit{Document} applicando un albero di filtri logici tramite 
    \textit{SearchDocumentsUC}, costruito a partire da un \textit{SearchRequestDTO} validato con Zod e trasformato in un 
    \textit{SearchDocumentsQuery} tramite \textit{MetadataKeyMapper}.
    \item \textbf{Ricerca semantica} – ricerca i \textit{Document} per similarità vettoriale tramite \textit{SearchSemanticUC}, 
    convertendo la query in un vettore di embedding tramite \textit{IWordEmbedding} e calcolando la \textit{cosine similarity} 
    rispetto ai vettori indicizzati.
    \item \textbf{Ricerca chiavi di metadati} – recupera le chiavi di metadati personalizzati presenti nei documenti del pacchetto ( quindi nel Database ).
\end{itemize}

Inoltre è esposto un canale per ottenere lo stato del motore AI, cioè se il modello è stato inizializzato e il numero di documenti indicizzati.

I risultati vengono mappati in DTO di presentazione (\textit{DocumentSearchResult}, \textit{ProcessSearchResult}, \textit{DocumentClassSearchResult}) 
prima di essere restituiti al frontend, mantenendo il dominio disaccoppiato dal livello di trasporto IPC.

\subsection{Architettura del modulo}

\begin{figure}[H]
    \centering
    \makebox[\textwidth][c]{\includegraphics[width=1.2\textwidth]{../../../Docs/out/3_PB/DocumentiEsterni/DiagrammiSpecificaTecnica/generated_/architecture/c1-c3/C3-main-search-flow/C3_Main_Search_Flow_No_Repos.png}}
    \caption{Diagramma C3 – Componenti del modulo di ricerca}
    \label{fig:C3_search_no_repos}
\end{figure}

\begin{figure}[H]
    \centering
    \includegraphics[width=0.8\textwidth]{../../../Docs/out/3_PB/DocumentiEsterni/DiagrammiSpecificaTecnica/generated_/architecture/c1-c3/C3-main-search-flow/C3_Main_Search_Flow_Repos.png}
    \caption{Diagramma C3 – Componenti del modulo di ricerca (Repository e DAO)}
    \label{fig:C3_search_repos}
\end{figure}

\subsubsection{SearchIpcAdapter}

\begin{figure}[H]
    \centering
    \includegraphics[width=1\textwidth]{../../../Docs/out/3_PB/DocumentiEsterni/DiagrammiSpecificaTecnica/C4/C4_backend_Search/C4_backend_Search.png}
    \caption{Diagramma delle classi – SearchIpcAdapter}
    \label{fig:search_ipc_adapter}
\end{figure}

\textit{SearchIpcAdapter} è l'adapter inbound che registra su \textit{IpcMain} i canali IPC del modulo di ricerca.
Risolve le dipendenze tramite il container DI (tsyringe) e instrada le richieste dal renderer process ai rispettivi Use Case, senza contenere logica di business.

\paragraph{Metodo principale}
\textit{+ register(ipcMain: IpcMain): void} – metodo statico che registra tutti i canali IPC del modulo. Per ciascun canale risolve dal container il 
corrispondente Use Case o adapter, ne invoca \textit{execute()} e restituisce i risultati mappati in DTO di presentazione.

\paragraph{Canali registrati}

\begin{itemize}
    \item \textit{SEARCH\_CLASSES} : ricerca le classi documentali per nome, delegando a \textit{ISearchDocumentalClassUC} e mappando i risultati 
    tramite \textit{DocumentClassMapper.toSearchResult()}.
    \item \textit{SEARCH\_PROCESSES} : ricerca i processi per UUID, delegando a \textit{ISearchProcessUC} e mappando tramite \textit{ProcessMapper.toSearchResult()}.
    \item \textit{SEARCH\_DOCUMENTS} : ricerca strutturata tramite filtri. Valida il \textit{SearchRequestDTO} con Zod, normalizza i path dei metadati 
    tramite \textit{MetadataKeyMapper.mapGroup()}, costruisce un \textit{SearchDocumentsQuery} e delega a \textit{ISearchDocumentsUC}. I risultati vengono
     mappati tramite \textit{DocumentMapper.toSearchResult()}.
    \item \textit{SEARCH\_SEMANTIC} : ricerca semantica testuale, delegando a \textit{ISearchSemanticUC} e mappando i risultati (con score di similarità) 
    tramite \textit{DocumentMapper.toSearchResult()}.
    \item \textit{SEARCH\_FULLTEXT} : ricerca full-text polimorfica che, in base al \textit{SearchQueryType} della query (\textit{PROCESS\_ID}, 
    \textit{CLASS\_NAME}, \textit{FREE}), instrada la richiesta al Use Case appropriato.
    \item \textit{SEARCH\_GET\_AI\_STATE} : restituisce lo stato del motore AI interrogando \textit{IWordEmbedding.isInitialized()} e 
    \textit{IDocumentRepository.getIndexedDocumentsCount()}, costruendo un oggetto con \textit{status}, \textit{progressPercentage} e \textit{indexedDocuments}.
    \item \textit{SEARCH\_CUSTOM\_METADATA\_KEYS} : recupera le chiavi di metadati personalizzati delegando a \textit{IGetCustomMetadataKeysUC}.
\end{itemize}

\subsubsection{Use Case di ricerca}

I cinque Use Case del modulo implementano ciascuno un'interfaccia dedicata e delegano le operazioni di persistenza ai rispettivi repository, senza contenere logica infrastrutturale.

\paragraph{SearchDocumentalClassUC}

\begin{figure}[H]
    \centering
    \includegraphics[width=0.9\textwidth]{../../../Docs/out/3_PB/DocumentiEsterni/DiagrammiSpecificaTecnica/C4/C4_backend_Search/C4_backend_Search-1.png}
    \caption{Diagramma delle classi – SearchDocumentalClassUC}
    \label{fig:uc_search_class}
\end{figure}

Implementa \textit{ISearchDocumentalClassUC}. Il metodo \textit{execute(name: string): DocumentClass[]} invoca\\
\textit{IDocumentClassRepository.searchDocumentalClasses(name)} e restituisce le entità corrispondenti, senza logica aggiuntiva.

\paragraph{SearchProcessUC}

\begin{figure}[H]
    \centering
    \includegraphics[width=0.9\textwidth]{../../../Docs/out/3_PB/DocumentiEsterni/DiagrammiSpecificaTecnica/C4/C4_backend_Search/C4_backend_Search-2.png}
    \caption{Diagramma delle classi – SearchProcessUC}
    \label{fig:uc_search_process}
\end{figure}

Implementa \textit{ISearchProcessUC}. Il metodo \textit{execute(uuid: string): Process[]} invoca \textit{IProcessRepository.searchProcesses(uuid)} 
e restituisce le entità corrispondenti.

\paragraph{SearchDocumentsUC}

\begin{figure}[H]
    \centering
    \includegraphics[width=0.9\textwidth]{../../../Docs/out/3_PB/DocumentiEsterni/DiagrammiSpecificaTecnica/C4/C4_backend_Search/C4_backend_Search-3.png}
    \caption{Diagramma delle classi – SearchDocumentsUC}
    \label{fig:uc_search_documents}
\end{figure}

Implementa \textit{ISearchDocumentsUC}. Riceve un \textit{SearchDocumentsQuery} già costruito dall'adapter. Il metodo 
\textit{execute(filters: SearchDocumentsQuery): Document[]} invoca \textit{IDocumentRepository.searchDocument(filters)} e restituisce i documenti 
corrispondenti ai criteri di filtro.

\paragraph{SearchSemanticUC}

\begin{figure}[H]
    \centering
    \includegraphics[width=0.9\textwidth]{../../../Docs/out/3_PB/DocumentiEsterni/DiagrammiSpecificaTecnica/C4/C4_backend_Search/C4_backend_Search-4.png}
    \caption{Diagramma delle classi – SearchSemanticUC}
    \label{fig:uc_search_semantic}
\end{figure}

Implementa \textit{ISearchSemanticUC}. Coordina \textit{IWordEmbedding} per la generazione del vettore di query e \textit{IDocumentRepository} per 
la ricerca tra i vettori indicizzati. Il metodo \textit{execute(query: string): Promise<SemanticSearchMatch[]>} converte la query in un vettore tramite 
\textit{IWordEmbedding.generateEmbedding()}, quindi invoca \textit{IDocumentRepository.searchDocumentSemantic()} e restituisce un array di 
\textit{SemanticSearchMatch}, ciascuno con il \textit{Document} e il relativo \textit{score} di similarità.

\paragraph{GetCustomMetadataKeysUC}

\begin{figure}[H]
    \centering
    \includegraphics[width=0.9\textwidth]{../../../Docs/out/3_PB/DocumentiEsterni/DiagrammiSpecificaTecnica/C4/C4_backend_Search/C4_backend_Search-5.png}
    \caption{Diagramma delle classi – GetCustomMetadataKeysUC}
    \label{fig:uc_get_metadata_keys}
\end{figure}

Implementa \textit{IGetCustomMetadataKeysUC}. Il metodo \textit{execute(dipId: number | null): string[]} invoca 
\\\textit{IDocumentRepository.getDistinctCustomMetadataKeys(dipId)} e restituisce le chiavi distinte. Se \textit{dipId}
 è \textit{null}, le chiavi vengono recuperate sull'intero database senza filtrare per pacchetto.

\subsubsection{Modelli di dominio per la ricerca}

\paragraph{SearchDocumentsQuery}

\begin{figure}[H]
    \centering
    \includegraphics[width=0.6\textwidth]{../../../Docs/out/3_PB/DocumentiEsterni/DiagrammiSpecificaTecnica/C4/C4_backend_Search/C4_backend_Search-10.png}
    \caption{Diagramma delle classi – Query Model: SearchDocumentsQuery, SearchGroup, SearchCondition}
    \label{fig:search_query_model}
\end{figure}

Oggetto query che incapsula l'albero di filtri logici da applicare alla ricerca dei documenti. Viene costruito dall'adapter a partire dal 
\textit{SearchRequestDTO} validato e passato allo Use Case e poi al DAO.

\begin{itemize}
    \item \textit{filter: SearchGroup} : radice dell'albero di filtri.
    \item \textit{isTooDeep(maxDepth: number): boolean} : calcola ricorsivamente la profondità dell'albero e restituisce \textit{true} 
    se supera \textit{maxDepth} (default 5), per prevenire query eccessivamente annidate.
\end{itemize}

\paragraph{SearchGroup}
Interfaccia che rappresenta un nodo dell'albero di filtri. Implementa il \textit{\textbf{Composite pattern}}: ogni elemento figlio può essere a sua 
volta un \textit{SearchGroup} (nodo) o una \textit{SearchCondition} (foglia).

\begin{itemize}
    \item \textit{logicOperator: LogicOperator} : operatore \textit{AND} o \textit{OR} che combina i figli.
    \item \textit{items: ReadonlyArray<SearchCondition | SearchGroup>} : elementi figli del gruppo.
\end{itemize}

\paragraph{SearchCondition}
Interfaccia che rappresenta una foglia dell'albero, ovvero una singola condizione atomica su un campo del documento.

\begin{itemize}
    \item \textit{path: string} : percorso dot-notation del campo (es.\ \textit{NomeDelDocumento} o\\\textit{CustomMetadata.codice}).
    \item \textit{operator: SearchOperator} : operatore di confronto: \textit{EQ}, \textit{GT}, \textit{LT}, \textit{LIKE}, \textit{IN}, \textit{ELEM\_MATCH}.
    \item \textit{value: any} : valore di confronto. Per \textit{IN} deve essere un array; per \textit{ELEM\_MATCH} un \textit{SearchGroup} annidato.
\end{itemize}

\paragraph{SearchRequestDTO e SearchRequestSchema}
DTO e schema Zod per la richiesta di ricerca strutturata dal frontend. \textit{SearchRequestSchema} valida ricorsivamente l'albero di filtri, garantendo la correttezza 
strutturale prima che la richiesta raggiunga il dominio. Il path di ogni \textit{SearchCondition} è validato tramite una regex che consente solo la dot-notation alfanumerica.

\paragraph{SearchQuery}
Interfaccia per le query di ricerca full-text o semantica provenienti dal frontend.

\begin{itemize}
    \item \textit{text: string} : testo inserito dall'utente.
    \item \textit{type: SearchQueryType} : tipo della query (\textit{PROCESS\_ID}, \textit{CLASS\_NAME}, \textit{FREE}), usato da \textit{SearchIpcAdapter} 
    per l'instradamento nel canale \textit{SEARCH\_FULLTEXT}.
    \item \textit{useSemanticSearch: boolean} : flag per attivare la ricerca semantica.
\end{itemize}

\subsubsection{Mapper di ricerca}

I mapper trasformano le entità di dominio nei DTO di presentazione restituiti al frontend. Sono metodi statici collocati nei rispettivi mapper di persistenza per 
coerenza architetturale.

\paragraph{DocumentClassMapper.toSearchResult()}

\begin{figure}[H]
    \centering
    \includegraphics[width=0.7\textwidth]{../../../Docs/out/3_PB/DocumentiEsterni/DiagrammiSpecificaTecnica/C4/C4_backend_Search/C4_backend_Search-12.png}
    \caption{Diagramma delle classi – DocumentClassMapper}
    \label{fig:document_class_mapper}
\end{figure}

Trasforma un'entità \textit{DocumentClass} in un \textit{DocumentClassSearchResult}, verificando che l'entità sia già persistita (id valido). Restituisce i campi 
\textit{id}, \textit{uuid}, \textit{integrityStatus}, \textit{name} e \textit{type} (fisso alle Classi).

\paragraph{ProcessMapper.toSearchResult()}

\begin{figure}[H]
    \centering
    \includegraphics[width=0.5\textwidth]{../../../Docs/out/3_PB/DocumentiEsterni/DiagrammiSpecificaTecnica/C4/C4_backend_Search/C4_backend_Search-13.png}
    \caption{Diagramma delle classi – ProcessMapper}
    \label{fig:process_mapper}
\end{figure}

Trasforma un'entità \textit{Process} in un \textit{ProcessSearchResult}, verificando la persistenza. Restituisce i campi \textit{id}, \textit{uuid}, 
\textit{integrityStatus} e \textit{type} (fisso ai processi).

\paragraph{DocumentMapper.toSearchResult()}

\begin{figure}[H]
    \centering
    \includegraphics[width=0.8\textwidth]{../../../Docs/out/3_PB/DocumentiEsterni/DiagrammiSpecificaTecnica/C4/C4_backend_Search/C4_backend_Search-14.png}
    \caption{Diagramma delle classi – DocumentMapper}
    \label{fig:document_mapper}
\end{figure}

Trasforma un'entità \textit{Document} e uno score opzionale in un \textit{DocumentSearchResult}. Estrae il nome cercando prima il nodo \textit{NomeDelDocumento}, 
poi \textit{nome}, nell'albero dei metadati. Determina il tipo (\textit{DocumentTypeEnum}) tramite \textit{mapRootMetadataType()}, che normalizza il nome del 
nodo radice in \textit{UPPER\_SNAKE\_CASE} e lo confronta con i valori canonici. Restituisce i campi \textit{id}, \textit{uuid}, \textit{integrityStatus}, 
\textit{name}, \textit{type} e \textit{score}.

\paragraph{MetadataKeyMapper}

\begin{figure}[H]
    \centering
    \includegraphics[width=0.8\textwidth]{../../../Docs/out/3_PB/DocumentiEsterni/DiagrammiSpecificaTecnica/C4/C4_backend_Search/C4_backend_Search-11.png}
    \caption{Diagramma delle classi – MetadataKeyMapper}
    \label{fig:metadata_key_mapper}
\end{figure}

Utility class che normalizza i path dei metadati dal frontend in PascalCase, garantendo coerenza con la struttura interna del JSON dei metadati.

\begin{itemize}
    \item \textit{toPascalCase(key: string): string} : converte una chiave in PascalCase, gestendo separatori \textit{\_} e \textit{-} e trasformando ogni segmento 
    di path dot-notation separatamente.
    \item \textit{mapGroup(group: SearchGroup): SearchGroup} : applica ricorsivamente \textit{toPascalCase()} al campo \textit{path} di ogni \textit{SearchCondition}
     nell'albero, inclusi i valori di condizioni \textit{ELEM\_MATCH} con \textit{SearchGroup} annidato.
    \item \textit{mapFilters(filters: MetadataFilter[]): MetadataFilter[]} : variante per array piatti di filtri chiave-valore; normalizza il campo \textit{key} di 
    ciascun filtro.
    \item \textit{fromLegacyFilters(filters: MetadataFilter[]): SearchGroup} : metodo di compatibilità che converte un array piatto di filtri legacy in un 
    \textit{SearchGroup} con operatore \textit{AND}, filtrando i valori nulli o vuoti e impostando l'operatore \textit{EQ} per ogni condizione.
\end{itemize}

\subsubsection{Repository e DAO di ricerca}
 
I repository implementano il \textit{\textbf{pattern Repository}} delegando le operazioni di persistenza ai rispettivi DAO, che interagiscono direttamente con 
SQLite tramite \textit{better-sqlite3}.
 
\paragraph{DocumentClassRepository e DocumentClassDAO}
 
\begin{figure}[H]
    \centering
    \includegraphics[width=0.7\textwidth]{../../../Docs/out/3_PB/DocumentiEsterni/DiagrammiSpecificaTecnica/C4/C4_backend_Search/C4_backend_Search-6.png}
    \caption{Diagramma delle classi – DocumentClassRepository e DAO}
    \label{fig:document_class_repo}
\end{figure}
 
\textit{DocumentClassRepository} espone \textit{searchDocumentalClasses(name)}, delegando a \textit{DocumentClassDAO.search(name)}, che esegue la query 
e mappa le righe risultanti in entità \textit{DocumentClass} tramite \\\textit{DocumentClassMapper.fromPersistence()}.
 
\paragraph{ProcessRepository e ProcessDAO}
 
\begin{figure}[H]
    \centering
    \includegraphics[width=0.7\textwidth]{../../../Docs/out/3_PB/DocumentiEsterni/DiagrammiSpecificaTecnica/C4/C4_backend_Search/C4_backend_Search-7.png}
    \caption{Diagramma delle classi – ProcessRepository e DAO}
    \label{fig:process_repo}
\end{figure}
 
\textit{ProcessRepository} espone \textit{searchProcesses(uuid)}, delegando a \textit{ProcessDAO.searchProcesses(uuid)}, che esegue la query e 
mappa le righe risultanti in entità \textit{Process}.
 
\paragraph{DocumentRepository e DocumentDAO}
 
\begin{figure}[H]
    \centering
    \includegraphics[width=\textwidth]{../../../Docs/out/3_PB/DocumentiEsterni/DiagrammiSpecificaTecnica/C4/C4_backend_Search/C4_backend_Search-8.png}
    \caption{Diagramma delle classi – DocumentRepository e DAO}
    \label{fig:document_repo}
\end{figure}
 
\textit{DocumentRepository} espone i seguenti metodi di ricerca:
 
\begin{itemize}
    \item \textit{searchDocument(filters: SearchDocumentsQuery): Document[]} : delega a \\\textit{DocumentDAO.searchDocument()}, che utilizza 
    \textit{DocumentMetadataQueryBuilder} per tradurre i filtri in una query ed eseguirla.
    \item \textit{searchDocumentSemantic(queryVector: Float32Array): Promise<Array<\{document, score\}> >} : delega a \textit{DocumentDAO.searchDocumentSemantic()},
     che esegue la ricerca per similarità vettoriale e restituisce le prime 10 coppie \textit{\{document, score\}} ordinate per score decrescente.
    \item \textit{getDistinctCustomMetadataKeys(dipId: number | null): string[]} : esegue la query e restituisce le chiavi di metadati personalizzati 
    distinte, deduplicate e ordinate alfabeticamente.
    \item \textit{getIndexedDocumentsCount(): number} : esegue la query e restituisce il numero di documenti con vettore indicizzato, o 0 se la tabella non esiste.
\end{itemize}

\subsubsection{DocumentMetadataQueryBuilder}
 
\begin{figure}[H]
    \centering
    \includegraphics[width=\textwidth]{../../../Docs/out/3_PB/DocumentiEsterni/DiagrammiSpecificaTecnica/C4/C4_backend_Search/C4_backend_Search-9.png}
    \caption{Diagramma delle classi – DocumentMetadataQueryBuilder}
    \label{fig:metadata_query_builder}
\end{figure}
 
\textit{DocumentMetadataQueryBuilder} traduce un \textit{SearchDocumentsQuery} in una clausola SQL parametrizzata compatibile con SQLite. Viene 
istanziato direttamente da \textit{DocumentDAO}.
 
\begin{itemize}
    \item \textit{build(filters: SearchDocumentsQuery): \{sql: string, params: unknown[]\}} : punto di ingresso pubblico. Verifica che il filtro non sia 
    vuoto, delega a \textit{buildQueryFromGroup()} sulla radice e restituisce la clausola SQL con i relativi parametri.
    \item \textit{buildQueryFromGroup(group: SearchGroup, jsonSource: string): \{sql, params\}} : costruisce ricorsivamente la clausola SQL per un 
    \textit{SearchGroup}, distinguendo nodi (ricorsione) e foglie (\textit{buildCondition()}). Le clausole vengono unite con \textit{AND} o 
    \textit{OR} in base al \textit{logicOperator}.
    \item \textit{buildCondition(condition: SearchCondition, jsonSource: string): \{sql, params\}} : traduce una singola \textit{SearchCondition} 
    nella clausola SQL corrispondente, gestendo separatamente l'operatore \textit{ELEM\_MATCH} (array annidati) e gli operatori scalari.
    \item \textit{buildScalarCondition(jsonExtract, operator, value): \{sql, params\}} : genera la clausola SQL per operatori scalari (\textit{EQ},
     \textit{GT}, \textit{LT}, \textit{LIKE}, \textit{IN}), eseguendo le opportune normalizzazioni di tipo e valore.
    \item \textit{toJsonPath(path: string): string} : converte un path dot-notation nella sintassi JSONPath di SQLite, validando ogni segmento e 
    sollevando un'eccezione in caso di segmenti non validi.
\end{itemize}